% Part: normal-modal-logic
% Chapter: syntax-and-semantics
% Section: normal-modal-logics

\documentclass[../../../include/open-logic-section]{subfiles}

\begin{document}

\olfileid[id]{nml}{syn}{nor}

\olsection{Logika Modal Normal}

Suatu logika modal dapat dipandang sebagai sebarang himpunan !!{formula}s
yang berperilaku cukup baik. Ada beberapa logika modal yang secara khusus
menarik, baik karena memiliki interpretasi dan penerapan yang menarik,
tetapi terutama karena secara teoretis telah dipahami dengan baik. Inilah
yang disebut logika modal normal.

\begin{defn}\ollabel{defn:modal-logic}
Suatu \emph{logika modal}~$\Log L$ adalah himpunan !!{formula}s dalam
bahasa modal dasar yang
\begin{enumerate}
\item memuat semua tautologi proposisional;
\item tertutup di bawah substitusi seragam; dan
\item tertutup di bawah modus ponens, yakni setiap kali $!A$ dan $(!A \lif
  !B) \in \Log L$, maka $!B \in \Log L$ pula.
\end{enumerate}
\end{defn}

\begin{ex}
Beberapa himpunan !!{formula}s yang memang tidak terlalu menarik, tetapi
memenuhi definisi ini dan karena itu terhitung sebagai logika modal, ialah:
\begin{enumerate}
\item himpunan semua tautologi;
\item himpunan semua !!{formula}s (logika modal tak konsisten).
\end{enumerate}
\end{ex}

\begin{defn}\label{defn:normal-modal-logic}
Suatu logika modal~$\Log L$ \emph{normal} jika dan hanya jika
\begin{enumerate}
\item logika itu memuat keduanya,
\begin{align}
\tag{K} \Box(p \lif q) & \lif (\Box p \lif \Box q) \\
\tag{Dual} \Diamond p & \liff \lnot \Box \lnot p,
\end{align}
\item dan tertutup di bawah nesesitasi, yakni setiap kali $!A \in
  \Log L$, maka $\Box !A \in \Log L$ pula.
\end{enumerate}
\end{defn}

Kita akan menjumpai dua cara utama untuk mendefinisikan logika modal. Cara
pertama bersifat teori-bukti, sebagai himpunan !!{formula}s yang dapat
diturunkan dalam sistem !!{derivation} tertentu. Cara kedua bersifat
teori-model, sebagai himpunan !!{formula}s yang valid dalam kelas model
tertentu. Ini mencerminkan keadaan dalam logika proposisional biasa (atau,
dalam hal ini, logika orde pertama). Di sana pun kita dapat memilih suatu
himpunan !!{formula}s dengan dua cara berbeda, misalnya sebagai himpunan
tautologi yang benar dalam semua penetapan nilai-kebenaran dan sebagai
himpunan semua !!{formula}s yang !!{derivable} dalam suatu sistem
!!{derivation}. Bagi logika modal normal, cara ganda untuk menentukan logika
modal ini dimungkinkan dengan menggunakan model relasional dan sistem bukti
aksiomatik (jenis Hilbert).
\end{document}
